\begin{claim}{D.2.45}{Proved}{The obstruction survives the fourteen-skeleton; the exact integral test}
Put $L=Y\langle7\rangle$ and use the quadratic coordinates
$\pi_{15}D_2L=\mathbb Z/8\{c\}\oplus\mathbb Z/2\{u\}$
of D.2.43.  Let $P_j$ be an actual finite-type CW skeleton of
$L$, with the specified nine-skeleton and no cells below seven.
The following statements hold.

The element $\partial(2c)$, which represents $h_Ml_4$,
is nonzero in every $P_j$ for $9\leq j\leq14$.
Its order there is either four or two.  Thus cells of dimensions
twelve through fourteen can at most reduce its order; only
fifteen-cells can first kill it.  Moreover the full-target
question is already determined on the actual fifteen-skeleton:
\[
 \pi_{14}P_{15}\xrightarrow{\ \cong\ }\pi_{14}L,
 \qquad
 \operatorname{im}JH_{P_{15}}=\operatorname{im}JH_L
             \quad\text{in }\pi_{15}D_2L.
 \tag{D.2.45a}
\]

For the full target, set
\[
 a\in H^7(L;\mathbb Z/4),\qquad
 x=\beta_4a\in H^8(L;\mathbb Z),\qquad b=\rho_4x,
 \tag{D.2.45b}
\]
where $a$ is the fundamental coefficient class and $\beta_4$
is the integral connecting homomorphism for multiplication
by four.  There is a well-defined coset
\[
 \theta\in H^{15}(L;\mathbb F_2)/
                      \operatorname{Sq}^1H^{14}(L;\mathbb F_2)
\]
such that, with $i:\mathbb F_2\to\mathbb Z/4$, $i(1)=2$,
\[
 a b=i\theta,\qquad \beta_{\mathbb Z}\theta=x^2.
 \tag{D.2.45c}
\]
Here $\beta_{\mathbb Z}$ is the integral first Bockstein.
If $JH(f)=2c$ for a stable map $f:S^{15}\to L$, then
\[
 \left\langle a b,\rho_4h_{\mathbb Z}(f)\right\rangle=2,
 \qquad
 \left\langle\theta,\rho_2h_{\mathbb Z}(f)\right\rangle=1.
 \tag{D.2.45d}
\]
The latter pairing is independent of the indicated ambiguity.
In particular $x^2=0$ would exclude $2c$ from the full
James--Hopf image. The square is subsequently evaluated in
D.2.50. For a specified sphere, (D.2.45d) does not fix its exact
James--Hopf value; D.2.51 strengthens the existential pairing
condition to an equivalence. D.2.4 remains open.
\end{claim}
\begin{proof}
\textbf{The integral transfer and its complete kernel.}
For a stable map between suspension spectra, the stable Hopf
invariant satisfies
\[
 \operatorname{tr}JH(f)
    =(f\wedge f)\Delta_{\mathrm{source}}-\Delta_{\mathrm{target}}f.
 \tag{D.2.45e}
\]
This is \cite[Theorem A(iii), p.~3, with proof pp.~12--13]
{Klein2026Hopf}; its Theorem B, p.~4, and full proof
pp.~13--16 identify the invariant with the Segal--Snaith
projection used in D.2.43.  The comparison checks weights
one and two and eliminates higher weights.  Thus the operation
in (D.2.45e) is the specified James--Hopf map, not an
unidentified operation with a similar obstruction property.
For a source sphere the reduced stable diagonal is zero,
so the right side is $-\Delta_L f$.

The integral label complex of the nine-skeleton has
$a_7,b_8,r_8,s_9$, with $db=4a$ and $ds=2r$.
In its ordinary tensor square the degree-fifteen cycles are
\[
 t=ab+ba,\qquad v=ar,\qquad w=ra,
\]
with relations $4t=2v=2w=0$.  Explicitly the degree-fifteen
boundary is $d(ab)=-4aa$, $d(ba)=4aa$; the degree-sixteen
boundaries are
\[
 \begin{aligned}
 d(bb)&=4(ab+ba),& d(as)&=-2ar,& d(sa)&=2ra,\\
 d(br)&=4ar,& d(rb)&=4ra,& d(rr)&=0.
 \end{aligned}
\]
Consequently
\[
 H_{15}(L\wedge L;\mathbb Z)
   =\mathbb Z/4\{t\}\oplus\mathbb Z/2\{v,w\}.
 \tag{D.2.45f}
\]
Higher target cells first contribute smash cells in dimension
seventeen, so this group and its named cycles are unchanged.

In the free $C_2$ resolution from D.2.43 the transfer is the
norm in resolution degree zero and is zero in positive
resolution degrees.  It is a chain map, since the norm
annihilates the first resolution differential.  Therefore
\[
 c=e_0ab+2e_1aa\longmapsto t,
 \qquad u=e_0ar\longmapsto v+w.
 \tag{D.2.45g}
\]
D.2.43 identifies the quadratic degree-fifteen homotopy group
with its integral homology.  Thus the composite of transfer
with integral Hurewicz has exactly the kernel
\[
       \ker\bigl(\pi_{15}D_2L\longrightarrow
                      H_{15}(L\wedge L;\mathbb Z)\bigr)
                    =\langle4c\rangle.
 \tag{D.2.45h}
\]
The independent finite certificate
\texttt{code/D45\_integral\_transfer.py} checks the tensor
boundaries, their integral quotient, and every element of
this transfer homomorphism.

\medskip\noindent
\textbf{The sharp remaining skeleton dimension.}
For a target $P_j$ of dimension at most fourteen,
$H_{15}(P_j;\mathbb Z)=0$.  Thus (D.2.45e), followed by
integral Hurewicz, gives
\[
             \operatorname{im}JH_{P_j}\subseteq\langle4c\rangle.
 \tag{D.2.45i}
\]
The quadratic group has already stabilized at the
nine-skeleton, as proved in D.2.43.  The same claim's
natural Goodwillie comparison applies to these six-connected
spaces and identifies $\ker\partial$ with
$\operatorname{im}JH$.  Since $2c\notin\langle4c\rangle$,
its actual boundary is nonzero.  Its original order is four,
so (D.2.45i) allows precisely orders four and two.  This proves
the assertion for all choices of intervening attaching maps;
no calculation of the individual twelve-cell compositions
is needed for nonvanishing.

Cells of dimension at least sixteen do not change
$\pi_{14}$, by cellular approximation applied to maps and
homotopies.  This proves the first part of (D.2.45a).
The cofiber spectrum of $P_{15}\to L$ has no cells below
sixteen.  Hence $\pi_{15}^sP_{15}\to\pi_{15}^sL$ is
surjective.  The map on the quadratic degree-fifteen groups
is an isomorphism, again by D.2.43.  Naturality of $JH$
then proves equality of images in (D.2.45a).  Sixteen-cells
can impose relations in its domain, but cannot add values
to its image.

\medskip\noindent
\textbf{The full-target coefficient pairing.}
The classes $a,b$ in (D.2.45b) are represented in the low
cellular model by the duals of $a_7,b_8$, respectively;
$b$ is zero on $r_8$.  Consequently $a\otimes b$ evaluates
as one on $t$ and zero on $v,w$.  If $JH(f)=2c$, equations
(D.2.45e)--(D.2.45g) give
\[
                  \Delta_*h_{\mathbb Z}(f)=-2t.
\]
Cup and diagonal compatibility gives the first equality
in (D.2.45d), since $-2=2$ modulo four. This pairing alone does
not specify the James--Hopf value: the transfer kernel is
$\langle4c\rangle$, and the quadratic group also has the separate
$u$ coordinate. D.2.51 uses all the mod-two products below to
remove the latter coordinate and proves an existential
sufficiency statement.

The required mod-two product vanishes on the actual $L$.
Here is the low Adams--Hilton check.  The degree-six and
degree-seven loop classes have images in the James tensor
algebra of $\Omega Y$ given by
\[
 u_6\longmapsto aa,\qquad D_7\longmapsto ab+ba,
                   \qquad E_7\longmapsto f.
\]
The four degree-thirteen products $uD$, $Du$, $uE$, $Eu$
have distinctive words $aaab$, $abaa$, $aaf$, $faa$,
respectively.  Their images are
linearly independent.  These are all possible quadratic
terms of degree thirteen in a minimal mod-two
Adams--Hilton model of a six-connected space; cubic terms
start in degree eighteen.  The quadratic differential,
dual to the reduced diagonal, must therefore be zero in
these four directions.  It follows that
\[
       H^7(L;\mathbb F_2)\,H^8(L;\mathbb F_2)=0.
 \tag{D.2.45j}
\]
This uses the actual model and comparison established in
D.2.26 and D.2.42, including its opened
\cite[\S3.1, p.~300]{FGHHalperinLemaire1995} interface.
It does not promote an integral word projection to an
integral diagonal computation.

The coefficient sequence
$0\to\mathbb F_2\xrightarrow{i}\mathbb Z/4\to\mathbb F_2\to0$
and (D.2.45j) now give a class $\theta$ with $ab=i\theta$,
unique modulo $\operatorname{Sq}^1H^{14}L$.
For its integral Bockstein, choose an integral cochain
$A$ representing $a$, with $dA=4B$; then $[B]=x$.
Choose an integral $C^{14}$ with
$dC\equiv A\smile B\pmod2$ and put
\[
 E^{15}=\frac{A\smile B-dC}{2}.
\]
This divides an even integral cochain.  Leibniz gives
$dE=2B\smile B$, so $\theta=[E\bmod2]$ and
$\beta_{\mathbb Z}\theta=x^2$, proving (D.2.45c).

For any integral cycle $z$ the same coefficient sequence gives
\[
                  \langle ab,\rho_4z\rangle
                      =2\langle\theta,\rho_2z\rangle.
 \tag{D.2.45k}
\]
Its value is unchanged on adding a first square: the
homological first Bockstein vanishes on the reduction of
an integral cycle.  Taking $z=h_{\mathbb Z}(f)$ proves
the second equality in (D.2.45d).

Finally $L$ is rationally acyclic and of finite type:
$Y=\Sigma^2K(\mathbb F_2,2)$ has these properties, and
its finite Postnikov base has only finite homotopy groups.
Thus its integral homology in each positive degree is finite.
If $x^2=0$, the coefficient exact sequence gives an
integral lift of $\theta$.  Integral cohomology pairs as
zero with every torsion integral homology class.  Equation
(D.2.45k) then excludes the required value.  Conversely,
the integral universal coefficient sequence shows that
vanishing of this pairing on \emph{all} integral
fifteen-cycles is equivalent to $\beta_{\mathbb Z}\theta=0$:
the classes annihilating reductions of integral cycles
are precisely the reductions of integral degree-fifteen
classes.  Indeed this annihilator is
$\operatorname{Ext}(H_{14}L,\mathbb F_2)$, and
$\operatorname{Ext}(H_{14}L,\mathbb Z)\to
\operatorname{Ext}(H_{14}L,\mathbb F_2)$ is surjective,
since $\operatorname{Ext}_{\mathbb Z}^2=0$.
This does not identify the smaller subgroup
of stable spherical cycles.

The ordinary Postnikov square does not determine the even
coefficient in (D.2.45k).  Its coefficient map here is
$\mathbb Z/4\to\mathbb Z/8$, $1\mapsto4$, and its
cochain is $A\smile dA=4A\smile B\bmod8$;
see \cite[\S2, formula (2.1), p.~899]{Thomas1959Pontrjagin}.
It kills the value two modulo four.  The suspension identity
of that source's Theorem I does not recover this lost
coefficient. The integral square is determined in D.2.50;
the pairing on the stable Hurewicz image remains unresolved.
\end{proof}
